\documentclass[12pt]{article}
\usepackage[utf8]{inputenc}
\usepackage[T1]{fontenc}
\usepackage[margin=1in]{geometry}
\usepackage{setspace}
\usepackage{xcolor}
\usepackage[colorlinks=true, urlcolor=blue!60!black, linkcolor=blue!60!black]{hyperref}
\setlength{\parindent}{2em}
\setlength{\parskip}{1.25em}
\renewcommand{\baselinestretch}{1.0}
\begin{document}
\begin{center}
{\LARGE \textit{Christian Legal Society v.\ Martinez}, 561 U.S.\ 661 (2010)\par}
\vspace{0.8em}
{\large Case Brief --- Four-Page Edition\par}
\end{center}
\vspace{1em}

\section{Memory Jogger}

A public law school may condition official student-organization recognition, with its funding and communication benefits, on an ``all-comers'' open-membership policy, because the policy is a reasonable, viewpoint-neutral condition on access to a limited public forum (\href{https://supreme.justia.com/cases/federal/us/561/661/}{561 U.S. 661}).

\section{Detailed Case Facts}

Hastings College of the Law ran a Registered Student Organization (RSO) program. Recognition brought eligibility for funds from a mandatory student-activity fee, access to official newsletters, bulletin boards, and e-mail, permission to reserve facilities, and use of the school's name and logo. RSOs had to follow Hastings' Nondiscrimination Policy, and in a 2005 Joint Stipulation the parties agreed that Hastings read it to require acceptance of all comers: any student could join or seek leadership regardless of status or beliefs.

The Hastings chapter of the Christian Legal Society (CLS) required voting members and officers to sign a Statement of Faith and excluded anyone engaged in ``unrepentant homosexual conduct'' or holding contrary beliefs, though anyone could attend meetings. Hastings denied CLS's 2004 application, making it the only group ever refused among more than 60, but allowed CLS to use rooms and general bulletin boards. CLS operated independently, holding weekly Bible studies and events.

\section{Procedural History}

CLS sued under \S~1983, alleging violations of free speech, expressive association, and free exercise; the group Hastings Outlaw intervened. The district court granted summary judgment to Hastings, and the Ninth Circuit affirmed in a two-sentence memorandum. The Supreme Court affirmed on June 28, 2010, remanding CLS's pretext claim if preserved.

\section{Judicial Votes}

Affirmed and remanded, \textbf{5--4}.
\begin{itemize}
\item \textbf{Majority:} Justice Ginsburg, joined by Justices Stevens, Kennedy, Breyer, and Sotomayor.
\item \textbf{Concurring:} Justice Stevens; Justice Kennedy.
\item \textbf{Dissenting:} Justice Alito, joined by Chief Justice Roberts and Justices Scalia and Thomas.
\end{itemize}

\section{Holding}

When a public university opens a recognition program as a limited public forum, free-speech and expressive-association challenges to membership conditions are analyzed together under limited-public-forum doctrine, not strict scrutiny. An all-comers requirement is reasonable in light of the forum's educational purposes and viewpoint neutral. Because the policy is neutral and generally applicable, \textit{Employment Division v. Smith} forecloses a free-exercise exemption. The holding reaches only a genuine all-comers policy; the Court did not decide the validity of the written Nondiscrimination Policy standing alone.

\section{Analysis of Opinions}

\textbf{Majority (Ginsburg).} The Court held CLS to its stipulation, rejecting its ``unseemly attempt to escape'' it. It merged the speech and association claims under forum analysis rather than the strict scrutiny of \textit{Roberts v.\ United States Jaycees} and \textit{Boy Scouts v.\ Dale}, because speech and association are intertwined here, strict scrutiny would erase a forum's defining limits, and Hastings offered a subsidy rather than a prohibition. Unlike \textit{Healy}, \textit{Widmar}, and \textit{Rosenberger}, no group was singled out for its views. Deferring to educators, the Court found the policy reasonable: it opens fee-funded leadership to every student, spares administrators from policing whether belief-based exclusions mask status discrimination, encourages tolerance, and tracks state law, and CLS kept ample alternative channels. The rule is ``textbook viewpoint neutral'' even if its incidental burdens fall unevenly (\textit{Ward v. Rock Against Racism}).

\textbf{Stevens (concurring).} Even the written policy is viewpoint neutral because it targets discriminatory conduct, not belief; a public school need not subsidize exclusionary groups.

\textbf{Kennedy (concurring).} The policy served the educational aim of engagement across differences, but CLS would have ``a substantial case on the merits'' if the policy were adopted to disadvantage a group's views or used by hostile students to infiltrate it.

\textbf{Alito (dissenting).} The record showed that Hastings rejected CLS under its written policy, while allowing secular groups to require ideological adherence, and announced the all-comers interpretation only later, suggesting pretext. \textit{Healy} controls: denying recognition substantially burdens association. The all-comers rule is unreasonable given the forum's purpose of viewpoint diversity among groups and exposes small religious groups to hostile takeovers.

\section{Examples: Future Applications}

\begin{enumerate}
\item \textbf{Same-Side (Upheld): A Uniform University Rule.} A university applies a written all-comers policy to all 120 clubs from the start and denies funding to a conservative club whose officers must endorse free-market principles. Uniformly applied, the policy is reasonable and viewpoint neutral.
\item \textbf{Same-Side (Upheld): A City Grant Program.} A city conditions community-center grants on open membership and denies them to an atheist club limited to nonbelievers, which may still rent rooms on standard terms. The city dangles a subsidy rather than prohibiting association.
\item \textbf{Opposite-Side (Invalid): Selective Enforcement.} A university deregisters a Christian group for requiring leaders to share its faith while letting secular and identity-based clubs require mission alignment, as in \textit{InterVarsity Christian Fellowship v.\ University of Iowa}. That is viewpoint discrimination outside \textit{Martinez}.
\item \textbf{Opposite-Side (Invalid): A Statewide Mandate.} A statute forces every private, off-campus Bible-study group to admit any resident as a member and officer, backed by fines. That is a prohibition, triggering strict scrutiny under \textit{Dale} and \textit{Jaycees}.
\item \textbf{Fence-Sitter: Hostile Takeover.} Under a genuine all-comers policy, fifteen opposed students join a seven-member Orthodox Jewish club and vote themselves into office. Kennedy's caveat supports the founders, but the majority left no test for infiltration claims.
\end{enumerate}

\section{Critique}

\textbf{Doctrinal and originalist critiques.} Critics say reliance on the stipulation let the Court decide an artificial case while ducking the common policy that bars religious criteria but permits secular ideological screening. An all-comers rule can undermine the forum's own purpose of diversity among groups, and the carrot-and-stick distinction sits uneasily with the unconstitutional-conditions doctrine and \textit{Rosenberger}. Scholars such as John Inazu and Michael McConnell fault the decision for subordinating the right of assembly and free exercise to judge-made forum tiers, and later cases such as \textit{Fulton} have cabined it.

\textbf{Progressive perspectives.} Many progressive scholars defend the result as affirming civic equality, since LGBTQ and minority-faith students need not subsidize groups that exclude them. Yet some warn that a sweeping all-comers rule threatens progressive affinity groups too, and that broad deference to administrators could weaken \textit{Healy v.\ James}.

\section{Key Quotations}

\begin{quote}
``The First Amendment shields CLS against state prohibition of the organization's expressive activity, however exclusionary that activity may be. But CLS enjoys no constitutional right to state subvention of its selectivity.''
\end{quote}
\noindent --- Justice Ginsburg, Opinion of the Court, 561 U.S. at 669.

\begin{quote}
``It is, after all, hard to imagine a more viewpoint-neutral policy than one requiring all student groups to accept all comers.''
\end{quote}
\noindent --- Justice Ginsburg, Opinion of the Court, 561 U.S. 661.

\begin{quote}
``Application of the less-restrictive limited-public-forum analysis better accounts for the fact that Hastings, through its RSO program, is dangling the carrot of subsidy, not wielding the stick of prohibition.''
\end{quote}
\noindent --- Justice Ginsburg, Opinion of the Court, 561 U.S. 661.

\end{document}
